% Spanish study reference — from chat (Canarias, verbs, weather, etc.)
% Compile: pdflatex study-from-chat.tex

\documentclass[10pt,landscape,a4paper]{article}
\usepackage[utf8]{inputenc}
\usepackage[T1]{fontenc}
\usepackage[left=1cm,right=1cm,top=1cm,bottom=1cm]{geometry}
\usepackage{multicol}
\usepackage{enumitem}
\usepackage{xcolor}
\definecolor{accent}{RGB}{0,51,153}
\setlist{nosep,leftmargin=*}
\setlength{\columnsep}{1em}
\raggedright

\begin{document}
\begin{center}
{\large\bfseries\color{accent} Spanish study reference}\\[0.2em]
\small From chat: Canarias, verbs, weather, por/para, articles, testimony, ordering, Mexico, etc.
\end{center}
\vspace{0.4em}

\begin{multicols}{2}

\section*{\color{accent}Ordering (caf\'e/bar)}
Hola, no hablo muy bien espa\~nol. \textbf{Quisiera} un tinto, una tortilla de patatas y un pedido de chistorra, por favor. (Quisiera = more polite than quiero.)

\section*{\color{accent}Testimony (relaci\'on con idiomas)}
1. \textquestiondown C\'omo te llamas y de d\'onde eres? Me llamo Jade, soy de Estados Unidos. 2. \textquestiondown D\'onde vives y d\'onde estudias espa\~nol? Vivo en Madrid, estudio espa\~nol en la escuela. 3. \textquestiondown Qu\'e lenguas hablas? Hablo chino e ingl\'es, aprendo espa\~nol. 4. \textquestiondown Por qu\'e estudias espa\~nol? Para hablar mejor y hacer amigos. 5. \textquestiondown Qu\'e haces para aprender? Leo libros, como comida espa\~nola, hablo con amigos. 6. \textquestiondown Qu\'e quieres hacer para mejorar? Hablar con gente que habla espa\~nol bien y viajar. 7. \textquestiondown Planes futuro? Quiero ser cient\'ifica de computadoras y trabajar en otros pa\'ises.

\section*{\color{accent}Canarias (info sheet template)}
\begin{itemize}[leftmargin=*,nosep]
\item La capital es Las Palmas y Santa Cruz (dos capitales).
\item La lengua oficial es el espa\~nol.
\item La moneda es el euro.
\item Tiene muchos habitantes (aprox. 2,2 millones).
\item El clima es bueno y c\'alido / subtropical.
\item Un producto importante es el pl\'atano.
\item Un plato t\'ipico es papas arrugadas con mojo.
\item Lugares de inter\'es: Parque Nacional del Teide, playas (Tenerife, Gran Canaria), Garajonay (La Gomera).
\end{itemize}

\section*{\color{accent}Verbs -- present}
\textbf{Ser:} soy, eres, es, somos, sois, son.\\
\textbf{Tener:} tengo, tienes, tiene, tenemos, ten\'eis, tienen.\\
\textbf{Hacer:} hago, haces, hace, hacemos, hac\'eis, hacen.\\
\textbf{-AR:} -o, -as, -a, -amos, -\'ais, -an (trabajar $\to$ trabajo, trabajas...).\\
\textbf{-ER:} -o, -es, -e, -emos, -\'eis, -en (leer $\to$ leo, lees...).\\
\textbf{-IR:} -o, -es, -e, -imos, -\'is, -en (escribir $\to$ escribo, escribes...).\\
\textbf{Person:} yo, t\'u, \'el/ella/usted, nosotros, vosotros, ellos/ustedes.

\section*{\color{accent}Articles}
el / la / los / las.\\
\textbf{Masc:} el museo, el libro, el trabajo, el vino, el mensaje, el idioma, el diccionario, \textbf{el baile}.\\
\textbf{Fem:} \textbf{la ciudad}, la comida, la lengua, la noche, la guitarra, la pel\'icula, la clase. (Not ``el ciudad''.)\\
Plural: los museos, las comidas, etc. \textbf{Ir a + art:} \textquestiondown Quieres ir \textbf{al} restaurante? (a + el = al). A la exposici\'on.

\section*{\color{accent}Por / para / porque}
\textbf{Para} + infinitive = purpose (para viajar, para ver series).\\
\textbf{Porque} + clause = reason (porque quiero aprender).\\
\textbf{Por} + noun = for/because of (por la comida).\\
Ejemplo: Estudio espa\~nol \textbf{para} ver series. Estudio \textbf{porque} quiero viajar \textbf{por} Espa\~na.

\section*{\color{accent}Weather (el tiempo)}
\textbf{hace} + calor / fr\'io / sol / viento.\\
\textbf{hay} + nubes.\\
\textbf{est\'a} + nublado / nevando / lloviendo.\\
\textbf{es} + (clima) tropical / templado / seco / h\'umedo / c\'alido.\\
En Santiago: est\'a nevando. En Bilbao: est\'a lloviendo. En C\'adiz: hace viento. En Sevilla: hace sol.

\section*{\color{accent}Question words}
\textbf{\textquestiondown D\'onde?} $\to$ lugar (where).\\
\textbf{\textquestiondown Cu\'antos/as?} $\to$ cantidad (how many).\\
\textbf{\textquestiondown C\'omo?} $\to$ modo (how).\\
\textbf{\textquestiondown Qu\'e?} $\to$ cosa. \textbf{\textquestiondown Cu\'al/es?} $\to$ cu\'al de varios.

\section*{\color{accent}Superlatives (mundo latino)}
Aconcagua (Argentina) = monta\~na m\'as alta de Am\'erica. Atacama (Chile) = m\'as seco. Arenal (Costa Rica) = volcanes m\'as activos. M\'exico = muy poblado; aguacates.\\
\textbf{Otros:} El Amazonas = r\'io m\'as largo del mundo. El chino = idioma m\'as hablado. Monte Elbr\'us = monta\~na m\'as alta de Europa. Oymyakon y Yakutsk = ciudades m\'as fr\'ias.\\
China / East Asia: clima variado, veranos calurosos y h\'umedos, inviernos fr\'ios y secos, monzones.

\section*{\color{accent}M\'exico (quiz)}
Tequila = una bebida. Enchiladas = un plato t\'ipico. Bandera = \'aguila (y serpiente). Moneda = peso. Dos ciudades m\'as pobladas = Ciudad de M\'exico y Guadalajara.

\section*{\color{accent}Useful vocab (from chat)}
withdraw from weed = dejar de consumir marihuana.\\
picaflor (ave) = hummingbird; (persona) = womanizer.\\
playboy = donju\'an / mujeriego.\\
fart = pedo. \textbf{esfuerzo} = effort (no ``pedo'').\\
I love Spain = Me encanta Espa\~na.\\
Spain = 西班牙 (Xībānyá). USA = 美国 (Měiguó). China = 中国 (Zhōngguó). Mexico = 墨西哥 (Mòxīgē).\\
``T\'ia, que va, no me da tiempo, voy ya a la Ressi pero tengo un 4\,\%'' $\to$ No way, I don't have time, I'm already going to the Ressi but I have 4\% battery.

\section*{\color{accent}Travel story (cuarto d\'ia)}
Aqu\'i todo es incre\'ible. Hay sol. Hoy estamos en Barcelona, que est\'a en Espa\~na. La gente es amable. La comida es deliciosa y el plato m\'as t\'ipico es la paella. Hace calor y el clima es agradable. Ma\~nana visitamos La Sagrada Familia. Dicen que es muy bonita. Despu\'es vamos a la playa.

\section*{\color{accent}Party script (principiante)}
\textquestiondown C\'omo est\'as? Estoy bien, \textquestiondown y t\'u? Un poco cansada pero feliz de estar aqu\'i. \textquestiondown Quieres pizza? S\'i, por favor. Tengo hambre. \textquestiondown Vamos a bailar despu\'es? \textquestiondown Perfecto! Vamos.

\section*{\color{accent}SMI (salario m\'inimo) -- key phrases}
El SMI es la cantidad m\'inima que un trabajador puede recibir, fijada por el Estado. En 2019 se elev\'o a 900 euros al mes (14 pagas). Objetivo: mejorar la vida de quienes ganan menos. Efectos: algunos contratos no renovados, menos nuevas contrataciones; muchos mantuvieron empleo y mejor salario. Teor\'ia: salario por encima del equilibrio puede reducir demanda de trabajo. En resumen: la subida buscaba mejorar ingresos de trabajadores con salarios bajos.

\section*{\color{accent}Historia (Gibraltar / CSCE)}
CSCE = Conference on Security and Cooperation in Europe (luego OSCE). Doc confidencial: Parsons escribe; habla con el Rey de Espa\~na. Tema: Gibraltar, di\'alogo UK--Espa\~na, Acuerdo de Lisboa, reuniones. Bilateral = entre dos partes. Espa\~na quer\'ia control del border; UK mantiene soberan\'ia; Gibraltar vot\'o quedarse brit\'anico. Resumen: territorio brit\'anico que Espa\~na reclama; border, pol\'itica, econom\'ia, identidad.

\end{multicols}
\vfill
\normalsize
\end{document}
